% CV CDI Data Engineer - Stephane KPOVIESSI
% Two-column layout adapted from the freelance CV.
% Compile target: pdfLaTeX

\documentclass[10pt,a4paper]{article}

\usepackage[a4paper,left=0.95cm,right=0.95cm,top=0.55cm,bottom=0.55cm]{geometry}
\usepackage[utf8]{inputenc}
\usepackage[T1]{fontenc}
\usepackage{lmodern}
\usepackage{xcolor}
\usepackage{enumitem}
\usepackage[hidelinks]{hyperref}
\usepackage{tabularx}
\usepackage{titlesec}
\usepackage{fontawesome5}

\definecolor{accent}{HTML}{1F4E79}
\definecolor{accentdark}{HTML}{163A5A}
\definecolor{text}{HTML}{2F3437}
\definecolor{muted}{HTML}{666D70}
\definecolor{rule}{HTML}{D5E2EC}
\definecolor{sidebg}{HTML}{F2F7FB}

\pagestyle{empty}
\setlength{\parindent}{0pt}
\setlength{\parskip}{3.6pt}
\setlist[itemize]{leftmargin=*, topsep=3pt, itemsep=3pt, parsep=0pt, partopsep=0pt}

\titleformat{\section}
  {\fontsize{11.2pt}{13pt}\selectfont\bfseries\color{accent}}
  {}{0pt}{}
  [\vspace{2pt}{\color{rule}\titlerule}\vspace{3pt}]
\titlespacing*{\section}{0pt}{10pt}{5pt}

\newcommand{\smallsection}[2]{%
  \vspace{8pt}
  {\fontsize{8.9pt}{10.45pt}\selectfont\bfseries\color{accentdark}%
  {\scriptsize\color{accent}#1}\hspace{3.5pt}#2}\par
  \vspace{3.2pt}
}

\newcommand{\contactline}[2]{%
  \noindent\makebox[0.34cm][l]{{\color{accent}\scriptsize#1}}%
  \begin{minipage}[t]{\dimexpr\linewidth-0.36cm\relax}%
  \fontsize{8.15pt}{9.5pt}\selectfont #2%
  \end{minipage}\par
  \vspace{1pt}
}

\newcommand{\sideitem}[2]{%
  \textbf{#1}\par
  #2\par
  \vspace{3.5pt}
}

\newcommand{\certitem}[2]{%
  \textbf{#1}\par
  {\fontsize{7.15pt}{8.1pt}\selectfont\color{muted}#2}\par
  \vspace{4.4pt}
}

\newcommand{\role}[3]{%
  \vspace{3pt}
  \begin{tabularx}{\linewidth}{@{}Xr@{}}
  \textbf{#1} & {\color{muted}#2}
  \end{tabularx}\vspace{0.8pt}\\
  {\color{muted}#3}\par
  \vspace{3pt}
}

\newcommand{\projectline}[2]{%
  \textbf{#1} - #2\par
}
\newcommand{\githublink}[2]{%
  \par{\footnotesize\color{accent}\faIcon{github}\ \href{#1}{#2}}\par
}

\newsavebox{\maincolumnbox}
\newlength{\maincolumnheight}

\begin{document}
\fontsize{10pt}{12.3pt}\selectfont
\sloppy

\begin{minipage}[t]{0.52\textwidth}
{\LARGE\bfseries\color{accentdark} Stéphane KPOVIESSI}\par
\vspace{1pt}
{\large\color{accent}Data Analyst | Data Engineer}\par
\vspace{1pt}
{\color{muted}Disponible en CDI}
\end{minipage}
\hfill
\begin{minipage}[t]{0.42\textwidth}
\raggedleft
\vspace{2pt}
\makebox[\linewidth][r]{{\normalsize\bfseries\color{accent}SQL \textbullet{} Python \textbullet{} Power BI \textbullet{} Databricks}}\par
\vspace{3pt}
\makebox[\linewidth][r]{{\small\color{muted}Performance commerciale, KPI, BI et qualité de données.}}
\end{minipage}

\vspace{6pt}
{\color{rule}\hrule height 1pt}
\vspace{7pt}

\begin{lrbox}{\maincolumnbox}
\begin{minipage}[t]{0.695\textwidth}
\vspace{0pt}
\raggedright

\section{Profil}
Ingénieur diplômé en Big Data \& Data Science, je développe des pipelines produisant des indicateurs métier fiables. Chez Allianz, j'ai modernisé des traitements SAS sur Azure Databricks, avec un gain supérieur à 8\texttimes{} sur le workflow mesuré, à capacité mémoire identique. J'associe Python/PySpark, SQL, contrôles qualité et reporting pour rendre les données exploitables par les équipes.

\section{Expériences}
\role{Allianz France - Ingénieur Data / Data Scientist}{Nov. 2025 - Avr. 2026}{Direction P4D, Advanced Data \& Climate - Paris}
\begin{itemize}
  \item Modernisation d'un datamart d'assurance : analyse de plus de 15 000 lignes SAS et migration vers Python/PySpark sur Azure Databricks, avec plus de 150 règles métier en architecture Médaillon.
  \item Accélération de plus de \textbf{8\texttimes{} du workflow mesuré} après migration de SAS vers Azure Databricks, à capacité mémoire identique.
  \item Fiabilisation des indicateurs métier par des contrôles qualité géospatiaux et des tests automatisés de parité SAS/Python et de non-régression.
\end{itemize}

\vspace{8pt}
\role{LYNE ET FRERES SARL - Data Analyst}{2023 - 2025}{Missions d'appui stratégique - Bénin / À distance}
\begin{itemize}
  \item Structuration de données commerciales, opérationnelles et agricoles avec Python/Pandas pour suivre les KPI d'activité, volumes et tendances.
  \item Développement de rapports et tableaux de bord Streamlit pour aider les équipes à interpréter les résultats et repérer les écarts.
  \item Automatisation des traitements récurrents et accompagnement des utilisateurs dans la lecture et l'exploitation des reportings.
\end{itemize}

\vspace{8pt}
\role{Groupe Sylvagreg - Machine Learning Engineer}{Mai 2024 - Août 2024}{Lille, France}
\begin{itemize}
  \item Développement d'une solution expérimentale de reconstruction 3D à partir d'images 2D, par photogrammétrie et modèles génératifs PyTorch.
  \item Automatisation du prétraitement et de l'inférence ; déploiement AWS/Docker et interface de visualisation pour tester et examiner les reconstructions.
\end{itemize}

\section{Projets BI \& Data}
\projectline{BI RETAIL SOLUTION}{PySpark, Airflow, PostgreSQL, Power BI, Docker}
Développement d'une chaîne de reporting ventes et supply chain : transformations PySpark orchestrées par Airflow, entrepôt PostgreSQL et tableaux de bord Power BI pour suivre les ventes et tendances produit.
\githublink{https://github.com/EngineerProjects/BI_Retail}{github.com/EngineerProjects/BI\_Retail}

\vspace{7pt}
\projectline{Seshat}{Python/Go, PostgreSQL/pgvector, RAG, Electron/React}
Conception d'un runtime open source d'agents Go et d'une chaîne de recherche documentaire pour automatiser les tâches et produire des réponses sourcées, avec contrôle des accès.
{\footnotesize\color{accent}\faIcon{globe}\ \href{https://seshat-ai.com}{seshat-ai.com}}\par

\vspace{7pt}
\projectline{ROBOT COLLECTEUR DE DÉCHETS}{AceObject - YOLOv8, OpenCV, Dijkstra}
Développement d'un pipeline YOLOv8/OpenCV pour détecter et suivre les déchets en les distinguant des obstacles, avec planification de collecte par Dijkstra. Exploration de l'apprentissage par renforcement pour la navigation.

\end{minipage}
\end{lrbox}
\setlength{\maincolumnheight}{\dimexpr\ht\maincolumnbox+\dp\maincolumnbox\relax}
\addtolength{\maincolumnheight}{45pt}

\begin{minipage}[t]{0.285\textwidth}
\vspace{0pt}

\hspace*{-0.18cm}%
\colorbox{sidebg}{%
  \hspace*{0.18cm}%
  \begin{minipage}[t][\maincolumnheight][t]{0.91\linewidth}
  \vspace{8pt}
  \fontsize{8.45pt}{10.35pt}\selectfont
  \raggedright

  \smallsection{\faIcon{address-card}}{Contact}
  \contactline{\faIcon{map-marker-alt}}{Cergy / Île-de-France}
  \contactline{\faIcon{phone}}{+33 7 48 59 47 13}
  \contactline{\faIcon{envelope}}{\href{mailto:oastephanekpoviessi@gmail.com}{oastephanekpoviessi@gmail.com}}
  \contactline{\faIcon{link}}{\href{https://linkedin.com/in/stephanekpoviessi}{linkedin.com/in/stephanekpoviessi}}
  \contactline{\faIcon{link}}{\href{https://kpoviessi-stephane.vercel.app}{kpoviessi-stephane.vercel.app}}
  \contactline{\faIcon{github}}{\href{https://github.com/EngineerProjects}{github.com/EngineerProjects}}

  \smallsection{\faIcon{layer-group}}{Compétences}
  \sideitem{Analyse \& BI}{KPI commerciaux, reporting, Power BI, Streamlit, Python/Pandas.}
  \sideitem{SQL \& données}{SQL, PostgreSQL, modélisation, qualité, contrôles, tests.}
  \sideitem{Pipelines \& cloud}{PySpark, ETL/ELT, Azure Databricks, Azure Data Lake, Airflow.}
  \sideitem{Industrialisation}{Docker, FastAPI, Git, CI/CD, automatisation de traitements.}

  \smallsection{\faIcon{graduation-cap}}{Formation}
  \sideitem{JUNIA ISEN Lille}{Diplôme d'Ingénieur - Big Data \& Data Science, 2026.}
  \sideitem{JUNIA ISEN Rabat}{Licence - Électronique \& Systèmes Numériques, 2023.}
  \sideitem{CPMS}{Maths, Physique et Sciences de l'ingénieur, 2020 - 2022.}

  \smallsection{\faIcon{certificate}}{Certifications}
  \certitem{\href{https://www.linkedin.com/learning/certificates/fed380ef1713eb9d1633ac7cd02b74bf58cf5f4bcf7e5fd47cb39bce19649bae}{Se préparer au métier d'analyste de données}}{Microsoft / LinkedIn Learning \textbullet{} 2024}
  \certitem{\href{https://www.freecodecamp.org/certification/amiche/data-analysis-with-python-v7}{Data Analysis with Python}}{freeCodeCamp \textbullet{} 2023}
  \certitem{\href{https://www.coursera.org/account/accomplishments/verify/LU2ZQ972BWNF}{Create ML Models in Microsoft Azure}}{Microsoft / Coursera \textbullet{} 2023}

  \smallsection{\faIcon{language}}{Langues}
  \textbf{Français} - langue maternelle\par
  \textbf{Anglais} - B2, Linguaskill Business\par

  \smallsection{\faIcon{users}}{{\fontsize{8.4pt}{10pt}\selectfont Engagements \& associations}}
  \textbf{JUNIA ALUMNI}\par
  {\color{muted}Commission Réseau}\par
  \vspace{3pt}
  \textbf{Vice-président}\par
  {\color{muted}Sept. 2024 - Août 2025}\par
  Animation du réseau alumni et événements professionnels.\par
  \vspace{4pt}
  \textbf{Responsable communication}\par
  {\color{muted}Sept. 2023 - Août 2024}\par
  Annonces et valorisation des initiatives.\par
  \vspace{3pt}
  \end{minipage}
}

\end{minipage}
\hfill
\usebox{\maincolumnbox}

\end{document}
